%======================================================================
\section{限界・未解決問題・Phase~2・3への展望}
\label{sec:limitations_and_outlook}

%----------------------------------------------------------------------
\subsection*{導入}
\label{subsec:s9_intro}

本節では、Phase~1（決定論的定式化；本論文）で構成した
ACS理論の理論的限界を体系的に開示し、未解決問題を整理し、
Phase~2およびPhase~3への研究方向を記述したうえで結論を述べる。
本節は新たなACS構成要件、三原理、TIPA機構、関係カテゴリー、
閉塞・崩壊条件または組織ACSの適用条件を追加する節ではない。
既節で与えた定義・命題・仮定の射程を変更せず、
それらが現在の枠組みのもとで何を主張していないかを明示する。

構成は、\S\ref{subsec:limitations_list}（理論的限界）、
\S\ref{subsec:phase1a}（Phase~1.$\alpha$ への委譲事項）、
\S\ref{subsec:open_problems}（未解決問題）、
\S\ref{subsec:phase_outlook}（Phase~2・3への展望）、
\S\ref{subsec:conclusion}（結論）である。

%----------------------------------------------------------------------
\subsection{理論的限界}
\label{subsec:limitations_list}

以下では、\S\ref{sec:acs_def}から\S\ref{sec:org_acs}までに
明示されたPhase~1の境界と、現在なお残る未確定部分を
L1--L18の連番で整理する。

\paragraph{L1：決定論的・離散時間・コンパクト空間への限定}
\label{lim:L1}
Phase~1は、決定論的・離散時間の状態遷移と、
コンパクトな状態・行動・内部モデル等の空間を中心とする枠組みである。
したがって、確率的遷移、確率分布、期待値、連続時間極限、
ジャンプ過程、非有界空間および可変次元状態空間を
一般理論としては扱わない。
この限定は、存在・閉性・コンパクト性等を扱う際の技術的明確性を与える一方、
現実の確率的・連続時間的・非有界な系への直接適用を保証しない。

\medskip
\noindent\textit{起源：\S\ref{sec:acs_def}および論文全体のPhase~1規約。
Phase~2：確率的遷移、連続時間化、局所コンパクト化その他の弱化を検討する。}

\paragraph{L2：外部環境状態の遷移・到達可能性（環境動学）の限定的な扱い}
\label{lim:L2}
\S\ref{sec:acs_def}では、ACSの存続に不可欠な環境要因を
必要環境空間$E$として与え、背景型として外部環境状態空間
$E_{\mathrm{ext}}=E^{\mathrm{bg}}\times O$を置く。
Phase~1では、当該ACSの行動に応答しうる行動応答成分の遷移を、
環境実現の側の写像$g_O$として与えるにとどめる
（定義~\ref{def:env_realization}）。
外生背景成分の列の生成則、環境状態間の到達可能性、
環境の確率的変動、他のACS個体との結合動学、
およびACSと環境の共進化は、
独立の環境動学としては定式化しない。
したがって、異なる環境状態の存在可能性と、
その状態へ実際に到達する動学は区別される。
自己--環境境界の再編成と外部環境状態空間の時間依存はL16で扱う。

\medskip
\noindent\textit{起源：\S\ref{sec:acs_def}の必要環境空間の構造的解釈と外部環境状態空間の分解。
Phase~2：背景成分の生成則、環境遷移の可到達性、確率的環境、他個体との結合動学を導入する。}

\paragraph{L3：認知的処理能力の未形式化}
\label{lim:L3}
Phase~1は認知的処理能力を形式的に定義せず、
単一の量としても導入しない。
\S\ref{sec:acs_def}の備考~\ref{rem:cognitive_capacity_policy}は、
量的容量、時間応答性、統合の質を区別すべきこと、
容量と実現量を区別すべきこと、
および能力の使い方に現れる偏向性とその程度（偏向強度）を能力に含めないことを
記述方針として述べるにとどまる。
したがって、認知的処理能力の高低を比較する指標、
三つの区分の測定、およびそれらと可行性・閉塞・崩壊との関係を、
Phase~1の結果からは論じられない。

\medskip
\noindent\textit{起源：\S\ref{sec:acs_def}、
備考~\ref{rem:cognitive_capacity_policy}。
Phase~1.$\alpha$／Phase~2：三つの区分の形式化
（時間応答性については所定の基準に到達するまでの段数を候補とする）と、
その測定を検討する。}

\paragraph{L4：同一個体性における非分岐性}
\label{lim:L4}
仮定~\ref{ass:no_branching}は、
時刻 $t$ のACS個体実現に対して、
同一個体の継続として対応づけられる時刻 $t+1$ の
ACS個体実現を高々一つとし、
コピー個体または複数の後続実現が生じる分岐を
Phase~1の対象から除外する。

したがって、完全複製、分岐的自己同一性、
複数の後続実現への同一性継承等を含む系については、
時間的自己同一性関係をそのまま適用することはできない。
組織ACSでは、組織の分割・分社のように、
同一組織の継続として対応づけられる後続の実現が複数になる場合がこれに当たる。
なお、この非分岐性は状態遷移写像 $f$ の一価性とは
別個の個体同一性上の仮定である。

\medskip
\noindent\textit{起源：\S\ref{sec:acs_def}、
仮定~\ref{ass:no_branching}、
定義~\ref{def:time_identity}。
Phase~2：個体実現間の時間的継続関係、
集合値対応または系譜構造等を候補として、
分岐を許す時間的自己同一性を再定式化する。}
\paragraph{L5：T-STABの閾値較正と推論射程}
\label{lim:L5}
\S\ref{sec:three_principles}のT-STAB
（定義~\ref{def:tstab}）は、
SM評価の下限、局所FRとEUの窓内充足、およびSM評価の一時刻変動を
有限時間窓上で統合する三原理時間統合評価条件である。
これは力学系論におけるLyapunov安定性、漸近安定性、
または構造安定性を意味するものではない。
下限参照値 $\underline{s}_{\mathrm{SM}}$、
変動上限 $\delta_{\max}$、
および窓幅 $W_T$（定義~\ref{def:tstab_params}）は
対象ACSの評価構造に相対する参照パラメータであり、
Phase~1では所与として扱われる。
その定量的較正則は一般には与えられていない。

また本論文は、T-STAB条件の違反を、
閉塞 $A_{i,t}=\emptyset$、崩壊 $\mathrm{Col}_{i,t}=1$、
負の関係品質を伴う環境結合、完全孤立、
またはPriorityの発火の十分条件として置かない。
これらへの接続を述べるには追加条件を要する。
したがってT-STABは、
Phase~1では三原理に関係する理論上の状態を有限窓で
記述・診断する条件であり、
一般的な収束、漸近安定性または長期予測を
単独で与えるものではない。

また、Phase~1は、FRの長期の側面について、
独立した無限時間判定条件を与えない。
定義~\ref{def:FR}は各時刻の局所FRを与え、
T-STABは有限窓内におけるその充足を評価する。
したがって本論文は、T-STABの成立から、
有限窓外における局所FRの持続、
または無限時間にわたる応答多様性の維持を一般的な帰結として置かない。
長期の側面の本格的な理論は、\S\ref{subsec:phase1a}に挙げる。

\medskip
\noindent\textit{起源：\S\ref{sec:three_principles}、
定義~\ref{def:FR}、備考~\ref{rem:sd_local_scope}、
定義~\ref{def:tstab_params}、定義~\ref{def:tstab}、
命題~\ref{prop:tstab_implies_principles}、
備考~\ref{rem:conv_tstab_scope}および備考~\ref{rem:tstab_descriptive_status}。
Phase~2：閾値・窓幅の較正則、パラメータ感度、
ロバスト性、およびT-STAB違反から各失敗状態への
追加条件付き接続を実装別に検討する。}

\paragraph{L6：TIPAの具体的実装写像と推定誤差の定量化}
\label{lim:L6}
\S\ref{sec:tipa}は、M1・M2・M3、PSLおよびISLの
機能的構造と最小要件を与えるが、各実装における
信号変換、評価、Priority生成、応答選択等の
具体的関数形を一意には固定しない。
また、M1の状態・環境推定 $\hat x_t$、$\hat e_t$ と
実現値 $x_t$、$e_t$ との誤差、
およびM2の代理評価と原理水準の真値との誤差について、
Phase~1は普遍的な誤差上界、収束性または一致を要求しない
（備考~\ref{rem:m1_estimate_nontruth}、備考~\ref{rem:error_bounds}）。
したがって、推定・評価の誤差が各実装でどの程度に収まるか、
その時間変化、および誤差が可行行動集合や閉塞到達へ与える影響は、
Phase~1の結果からは定まらない。
TIPAはPhase~1では構成的・機能的形式を与えるが、
特定の生物学的機構またはAI実装を一意に同定する理論ではない。
また、構造的分離 $B_{SE}$ とM1-bの自己・環境分類の不整合を検証する機構
（備考~\ref{rem:xi_se_interpretation}）、
同一の $f$ に対応しうる内部モデル構造 $M$ の一意性
（備考~\ref{rem:M_independence}）、
内部表現の空間 $\mathcal{R}_{\mathrm{int}}$ から $\mathcal{M}$ への符号化写像
（備考~\ref{rem:M_encodes_representation}）、
および $f_{\mathrm{auto}}$、$f_{\mathrm{dir}}$、$\oplus$ の具体的実装形式
（仮定~\ref{ass:f_decomposition}）も定めていない。

\medskip
\noindent\textit{起源：\S\ref{sec:tipa}、
備考~\ref{rem:m1_estimate_nontruth}、備考~\ref{rem:error_bounds}。
Phase~2：実装別写像の具体化、誤差上界を導入するための条件、
推定誤差の動学、ロバスト性および実証的較正を検討する。
整合性検証機構、$M$ の一意性および $f$ の分解の実装形式も
Phase~2以降の課題とする。
符号化写像は、後続Phaseで導入することを約束しない。}

\paragraph{L7：\S\ref{sec:extended_dynamics}更新写像の抽象性と変数間結合の未確定}
\label{lim:L7}
\S\ref{sec:extended_dynamics}は、
$\boldsymbol{\beta}_t$、$d_{\mathrm{eff},t}$、$\pi_{\tau,t}$、
$\widehat{\tau}_{t}$および四健全度等について、
固定した読み出し規則により、状態・内部モデル・外部環境状態、
または作動記録の有限区間から定まる値として扱い、
その時刻$t$から$t+1$への変化は、$f$、Update、および外部環境の実現（背景成分の列と $g_O$）を通じて定まる。
独立した更新写像は置かない。
一方、具体的学習率、回復率、収束速度、変数間の相互依存関数、
長期安定性および一般的な閉形式解は固定しない。
また、$\widehat{\tau}_{t}=\pi_{\tau,t}\tau$は派生量であり、
独立した更新則をもたない。
したがって、\S\ref{sec:extended_dynamics}は時間方向の最小構造を与えるが、
特定実装の詳細な学習・劣化・回復モデルを与えるものではない。

\medskip
\noindent\textit{起源：\S\ref{sec:extended_dynamics}。
Phase~2：具体的更新則（予測と実現の差にもとづく内部モデルの更新の規則を含む）、結合構造、収束・安定性・回復条件をモデル別に検討する。
偏向性については、各偏向性の参照基準と尺度の具体化、偏向強度の読み出し方式、TIPA・PSL・ISL による形成の記述、形成の組み合わせと帰属区分の対応、観測者モデルに基づく帰属区分の同定、実装上の偏向性の変容・統合と抽象的構造記述の添字との対応づけ、および各偏向性の方向・内容の表現を検討する
（\S\ref{sec:extended_dynamics}、定義~\ref{def:bias_attribution}、備考~\ref{rem:beta_multiple_biases}）。}

\paragraph{L8：関係記述子の非因果性と測定・分類の実装依存性}
\label{lim:L8}
\S\ref{sec:rel_min}の関係記述子は、
対象間の関係構造と、関係継続窓において観測される
可行性関連量の変化を記述する分類Cの分析形式である。
有限窓差分
$\Delta_L\mathrm{sz}_i(A_{i,t})$、
$\Delta_Ld_{\mathrm{eff},i,t}$、
$\Delta_L\widehat{\tau}_{i,t}$は、
対象$Y_j$の一意な因果的寄与を意味しない。
また、局所的サイズ指標$\mathrm{sz}_i$の具体的構成、
カテゴリー判定閾値、持続時間、複合カテゴリー、因果帰属等は
一般理論として固定していない。
局所的サイズ指標の測度論的な精密化
（備考~\ref{rem:finite_window_diff_limit}）も行わない。

\medskip
\noindent\textit{起源：\S\ref{sec:rel_min}。
Phase~1.$\alpha$：分類・測定規則の詳細化。
Phase~2：介入・反実仮想等を含む因果的拡張を必要に応じて検討し、
局所的サイズ指標の測度論的な精密化を扱う。}

\paragraph{L9：閉塞・崩壊の観測可能性、回復判定および経路動学}
\label{lim:L9}
\S\ref{sec:collapse_cond}では、
閉塞$A_{i,t}=\emptyset$と
崩壊$\mathrm{Col}_{i,t}=1$を明確に分離し、
閉塞経路集合・崩壊到達経路集合・到達時刻を定義した。
しかし、有限観測だけから崩壊を一般に判定できるとは主張せず、
外的回復可能性についても実装依存の決定論的判定規則を抽象化している。
さらに、経路集合への流入率、到達確率、平均到達時間、
回復確率等はPhase~1では定義しない。
整形式実現指標 $\mathrm{WF}_{i,t}$ の値を有限観測から判定する操作的手続き
（定義~\ref{def:wf_realization}）と、
条件 (FT-1)・(FT-2)・(SC-1)・(SC-2) の判定規則の一般形式
（定義~\ref{def:functional_termination}、定義~\ref{def:collapse}）も与えていない。

\medskip
\noindent\textit{起源：\S\ref{sec:collapse_cond}。
Phase~2：観測可能性、判定アルゴリズム、確率的到達・回復解析、回復余裕の定量化を検討する。
$\mathrm{WF}_{i,t}$ の操作的判定手続きもPhase~2の課題とする。
(FT-1)・(FT-2)・(SC-1)・(SC-2) の判定規則は、
Phase~1では対象系ごとに固定されるものとして扱う。}

\paragraph{L10：参加ACS集合の時刻別記録と組織同一性動学の未形式化}
\label{lim:L10}
\S\ref{sec:org_acs}では、
組織 $\mathcal O$ の参加ACS集合
$\mathbb{I}_{\mathrm{org},t}(\mathcal O)$ と参加ACS数
$n_{\mathrm{org},t}(\mathcal O)$ を時刻ごとに記録する。
固定した有限候補集合 $P$ は、採用するモデルで必要な場合に限る
追加的な記述であり、Phase~1の一般要件ではない。
したがって、Phase~1は構成員数を時間的に固定する理論ではない。

一方、参加ACS集合そのものの一般的な遷移則、
加入、離脱、生成、消滅等の動学、
およびそれらの変化を通じた組織同一性の保持条件は定式化しない。
また本論文は、参加ACS数の変化のみから、
組織ACS資格の成立・不成立、組織の崩壊、
または同一組織ACSの復帰を一般的な帰結として置かない。

\medskip
\noindent\textit{起源：\S\ref{sec:org_acs}、
定義~\ref{def:org_participation}、
備考~\ref{rem:org_qualification_loss}および\S\ref{subsec:org_scope}。
Phase~2：参加集合の遷移、加入・離脱、
組織同一性および復帰条件の動学を対象モデルごとに検討する。}

\paragraph{L11：組織水準7タプルの存在・構成・一価遷移は一般定理ではない}
\label{lim:L11}
\S\ref{sec:org_acs}の
組織水準のACS構造
$\Sigma_{\mathrm{org}}$ は、
組織自身について
\S\ref{sec:acs_def}のACS定義、構成要件および整形式性が
満たされる場合に用いる構造形式である。
任意の社会的組織について、
$X_{\mathrm{org}}$、$U_{\mathrm{org}}$、
$E_{\mathrm{org}}$、$f_{\mathrm{org}}$、
$\Pi_{\mathrm{org}}$、$M_{\mathrm{org}}$、
$\mathrm{Update}_{\mathrm{org}}$ が
自動的に整形式に構成できることを証明するものではない。

また、組織状態は個体状態の和、平均または直積として
一般的に定まるものではなく、
集約写像を用いる場合には
同一の組織状態・組織行動・外部環境状態から
一価の組織次状態が定まることを別途確認する必要がある。
組織ACSの適用規約も、一般的存在または
任意の実在組織への該当性を主張しない。

\medskip
\noindent\textit{起源：\S\ref{sec:org_acs}、
定義~\ref{def:org_structural_form}、
備考~\ref{rem:org_state_closure}および定義~\ref{def:org_acs}。
Phase~2：具体的組織モデルに対する存在条件、
整形式性、一価遷移、識別可能性および実証的判定を検討する。}

\paragraph{L12：組織水準の判定規則と個体--組織集約則}
\label{lim:L12}
\S\ref{sec:org_acs}では、
組織自身に整形式なACS構造が成立する場合に限って、
先行節の三原理、閉塞・崩壊、時間統合評価および経路集合の定義を
組織水準へ適用する（備考~\ref{rem:org_principles_conditional}）。
その適用には、条件 (SC-1)・(SC-2)・(FT-1)・(FT-2) の判定規則を
組織について固定することを要する。
その一般形式は、個体の場合と同じく与えていない。

個体の可行性、三原理の時刻別充足述語の真偽、関係記述、健全度等を
平均・最小・加重平均その他の作用素によって
組織量へ変換する普遍的な集約則は与えない。
また、組織ACS資格、可行性、閉塞、
整形式実現の不成立および崩壊は
相互に同一視されない。
組織の分割・分社のように後続の実現が複数になる場合は、
非分岐性の仮定（L4）のもとでは扱えない。

\medskip
\noindent\textit{起源：\S\ref{sec:org_acs}、
備考~\ref{rem:org_individual_nonimplication}、
備考~\ref{rem:org_principles_conditional}および
備考~\ref{rem:org_failure_scope}。
Phase~2：個体--組織間の集約則、
組織水準の三原理評価、時間統合、
閉塞・崩壊経路および組織動態を個別モデルとして検討する。}

\paragraph{L13：ドメイン横断的対応と実証的妥当性の未検証}
\label{lim:L13}
本論文は、人間・生物、AI・ロボット、組織等へ
適用可能な抽象構造を意図しているが、
個別ドメインにおける7タプル各成分、三原理、TIPA、PSL・ISL、
関係記述子、閉塞・崩壊条件等の測定可能性と実証的妥当性を
本論文内で検証していない。
したがって、形式的適用可能性と経験的妥当性を区別する必要がある。
既存理論との関係についても、本論文が主張するのは
明示した範囲での整合性であり、包含・優越・一般化を意味しない。
TIPA成分の更新の質を表す量（定義~\ref{def:update_quality}）の
実証的な測定・推定方法も与えていない。

\medskip
\noindent\textit{起源：\S\ref{sec:intro}および全節の適用範囲。
Phase~2以降：ドメイン別操作化、測定、比較研究および反証可能性の検討を行う。}

\paragraph{L14：ECITフルモデルおよび規範理論はPhase~1の外部にある}
\label{lim:L14}
本論文は、倫理を善悪判断そのものとして定義するのではなく、
可行性保存に関する構造的制約として記述する。
したがって、愛・智慧・責任等のECIT諸要素、
倫理的協調応答の完全な発動条件、社会的規範の生成・変化、
具体的価値選択および政策的最適化をPhase~1の形式的帰結としては主張しない。
これらはACSの記述理論とは論理的地位を分けて扱う必要がある。

\medskip
\noindent\textit{起源：\S\ref{sec:intro}および\S\ref{sec:org_acs}の
記述理論・適用範囲としての限定。
Phase~3：ECITとの正式接続を、
Phase~1・2の結果と区別して検討する。}

\paragraph{L15：行動の実行の成否の区別}
\label{lim:L15}
\S\ref{sec:acs_def}の状態遷移写像の読みの規約
（備考~\ref{rem:f_reading}）のもとでは、
行動の試行、行動を実行できなかった場合、
および実行の成功の度合いを、$f$ の記述の中で区別しない。
行動入力の実行可能性を $U$ とは別に表す構造も置かないため、
$U$ の元をある時点で実行できないことは
$f$ が与える次状態に吸収され、
その行動が当該ACSの行動空間に含まれないこととは
別の形式対象として表されない。
また、作用差（定義~\ref{def:action_difference}）の大きさ、
多段階にわたる作用、行動や行動応答成分の寄与の比較、
および外部と内部に共通する資源量の収支も定量化しない。

\medskip
\noindent\textit{起源：\S\ref{sec:acs_def}、備考~\ref{rem:f_reading}。
Phase~2：実行の成否または程度を表す変数、
および時点ごとに実行可能な行動を $U$ とは別に表す構造を検討する。
作用差の大きさ等の定量化は、後続Phaseへの委譲事項または
将来の拡張・再検討の候補として保留する（\S\ref{sec:acs_def} の末尾の列挙）。}

\paragraph{L16：個体の構造の時間変化}
\label{lim:L16}
Phase~1は、個体の7タプル $\Sigma_i$ を時刻に依存しない構造体として扱い、
時刻とともに変わるのは状態、内部モデル状態その他の時刻別の値に限る
（\S\ref{sec:acs_def}、\S\ref{subsec:acs_7tuple} の記述水準の規約、定義~\ref{def:s_struct}）。
したがって、行動空間 $U$ そのものの拡大（新たな行動の獲得）または縮小、
内部モデルの機構構成の変化、
および状態空間の変化を、
同一個体の時間発展として記述しない。
可行集合を定める閾値の時間依存（定義~\ref{def:Pi}）、
行動空間 $U$ のラベル構造が位相構造に与える影響
（備考~\ref{rem:action_space_structure}）も扱わない。
また、7タプルの成分ではないが、
自己--環境境界の動的再編成（備考~\ref{rem:TSC_character}）、
時間依存する外部環境状態空間（仮定~\ref{ass:Eext_time_fixed}）、
および個体別の外部環境状態空間を相互に対応づける構造
（備考~\ref{rem:Eext_relative}）も扱わない。
他のACS個体との結合動学はL2に含める。

\medskip
\noindent\textit{起源：\S\ref{sec:acs_def}、\S\ref{subsec:acs_7tuple} の記述水準の規約、
備考~\ref{rem:TSC_character}、備考~\ref{rem:Eext_relative}、
備考~\ref{rem:action_space_structure}、仮定~\ref{ass:Eext_time_fixed}、定義~\ref{def:Pi}。
Phase~2：時刻に依存する個体の構造と、
その変化のもとでの時間的自己同一性を検討する。
動的閾値への拡張はPhase~2以降の課題とする。}

\paragraph{L17：必要環境要因の操作的同定・分解・判定}
\label{lim:L17}
\S\ref{sec:acs_def}は、必要環境要因を充足判定写像 $\eta\in E$ として与え、
構成要件CR-1（\S\ref{subsec:environment_space}）によってその数学的地位を定める。
一方、実際の対象から要因を同定し、その充足を判定する手続きは与えない。
具体的には、
必要環境要因の経験的同定手続きと同定誤差の扱い
（備考~\ref{rem:necessity_identification}）、
外部依存が一部の内部状態においてのみ現れる要因を要因と認める範囲とその程度、
CR-1の充足を有限観測から判定する手続き、
要因の同定が外部背景の取り方に依存する程度
（備考~\ref{rem:cr1_relativity}）、
要因の分解の一意性を与える追加条件
（備考~\ref{rem:factor_nonuniqueness}）、
および要因の分類粒度（備考~\ref{rem:eta_coarse_classes}）を定めていない。
必要性の存立と同定可能性は区別される。

\medskip
\noindent\textit{起源：\S\ref{sec:acs_def}、\S\ref{subsec:environment_space}、
備考~\ref{rem:necessity_identification}、備考~\ref{rem:cr1_relativity}、
備考~\ref{rem:factor_nonuniqueness}、備考~\ref{rem:eta_coarse_classes}。
Phase~2：経験的同定手続き、同定誤差、外部依存の範囲と程度、
有限観測によるCR-1の判定、外部背景への依存、分解の一意性の条件を検討する。
分類粒度の確定は、後続Phaseへの委譲事項または
将来の拡張・再検討の候補として保留する。}

\paragraph{L18：三原理に関わる補助量と解析条件の精密化}
\label{lim:L18}
\S\ref{sec:three_principles}は、三原理の時刻別充足述語と
その条件付き解析に必要な範囲で補助量を定める。
次の精密化は行わない。
環境応答成分 $I_E$ の同定に関する、観測ノイズ・時間遅延・内部フィードバック等を考慮した
操作的・因果的な決定基準（備考~\ref{rem:env_response_observation}）、
境界点における $x_E$ 方向の局所感応性の方向解析（備考~\ref{rem:pi_E_local_sensitivity}）、
$\Pi$ を補完する診断指標の動学的利用（備考~\ref{rem:diagnostic_role}）、
補題~\ref{lem:finite_exit} の履歴依存方策への一般化と離脱時刻の一様な有界性の条件
（備考~\ref{rem:uniform_bound}）、
応答多様度の値と閉塞到達時間との定量的関係（備考~\ref{rem:sd_degree}）、
$\mu_{\mathrm{eff}}$ の測度論的な精密化（備考~\ref{rem:mu_measure}）、
および仮定~\ref{ass:theta_exterior_blockage} を内部動学から導くための十分条件である。

\medskip
\noindent\textit{起源：\S\ref{sec:three_principles}、
備考~\ref{rem:env_response_observation}、備考~\ref{rem:pi_E_local_sensitivity}、
備考~\ref{rem:diagnostic_role}、備考~\ref{rem:uniform_bound}、
備考~\ref{rem:sd_degree}、備考~\ref{rem:mu_measure}、
仮定~\ref{ass:theta_exterior_blockage}。
Phase~2：上記の各項目を検討する。}

%----------------------------------------------------------------------
\subsection{Phase~1.$\alpha$ への委譲事項}
\label{subsec:phase1a}

本論文でいうPhase~1.$\alpha$は、
Phase~1の前提（決定論的・離散時間・コンパクト空間・非分岐性・単一ACS中心）と
各対象の型を変えずに、
Phase~1で値域・型・担体だけを与えた対象について、
規則・閾値・操作的基準を具体化する補助的詳細化の段階をいう。
Phase~1の仮定を弱める拡張はPhase~2に、
ECITとの接続はPhase~3に属する
（\S\ref{subsec:phase_outlook}）。
本文で「Phase~1.$\alpha$またはPhase~2」とした事項は、
決定論的に扱える部分をPhase~1.$\alpha$で、
それ以外の部分をPhase~2で検討する。
以下は検討課題の所在を示すものであり、
各事項がPhase~1.$\alpha$で形式化されることを主張しない。

\begin{enumerate}
  \item \textbf{構成的必要機能の候補：}
  必要環境利用機能（備考~\ref{rem:env_utilization_function}）、
  環境変化への横断的応答能力（備考~\ref{rem:env_change_response}）。
  \item \textbf{機構の構成：}
  TIPAの機構統合を許す再定式化（観察~\ref{obs:tipa_integration}）、
  保護的な働きの定まらない不随意の局所応答を含む、PSL層の局所自動応答一般
  （\S\ref{subsubsec:psl4} のPSL-4の範囲の外にあるもの）。
  \item \textbf{認知的処理能力：}
  量的容量・時間応答性・統合の質の形式化（備考~\ref{rem:cognitive_capacity_policy}）。
  \item \textbf{三原理：}
  行動応答成分の応答を考慮した応答不能環境集合・応答多様度・条件付き解析の版
  （備考~\ref{rem:same_time_alternatives}）、
  FRの長期の側面（備考~\ref{rem:sd_local_scope}）、
  \S\ref{sec:acs_def}の仮定と構成要件を満たすモデルによる
  時刻別充足述語の非含意の構成（備考~\ref{rem:principle_implication_scope}）。
  \item \textbf{TIPAの内容：}
  要請の質の列挙と分類（定義~\ref{def:request_signal_types}）、
  手段と目的の関係（備考~\ref{rem:purpose_hierarchy}）、
  目的を達成せずに終える理由の区別（定義~\ref{def:purpose_status}）、
  評価機構の迂回・置換の由来の追跡（備考~\ref{rem:self_attribution_scope}）、
  内部の対象トークンと外部の枠・個物の対応（備考~\ref{rem:Eext_decomposition_reading}）。
  \item \textbf{時刻間の量と記録：}
  実効的発展自由度の算出規則（定義~\ref{def:deff_section5}）、
  偏向性の形成の機構別の記述と帰属区分の同定（備考~\ref{rem:bias_formation_vs_attribution}）、
  保持記録が一致するための $\mathrm{Update}$ の詳細型の条件
  （定義~\ref{def:s5_record_agreement}、備考~\ref{rem:update_type_two_layers}）。
  \item \textbf{関係性：}
  \S\ref{sec:rel_min}の備考~\ref{rem:phase1a_delegation}の各項、
  局所的サイズ指標の具体化と有限近似（備考~\ref{rem:finite_window_diff_limit}）、
  主要方向の定量的判定規則（備考~\ref{rem:dir_undefined}）。
  \item \textbf{組織ACS：}
  組織ACSを関係対象とする記述と双方向合成の適用（備考~\ref{rem:org_terms_nonimplication}）。
\end{enumerate}

%----------------------------------------------------------------------
\subsection{未解決問題}
\label{subsec:open_problems}

以下はPhase~1では解決していない研究問題である。
各問題は将来の研究方向を示すものであり、
解決可能性、唯一の解法または必然的な進展を主張しない。

\begin{problem}[分岐を許す場合の同一個体性の再定式化]
\label{prob:branching_identity}

\textbf{背景：}
\S\ref{sec:acs_def} の仮定~\ref{ass:no_branching}
（非分岐性；強仮定）は、
時刻 $t$ のACS個体実現に対して、
同一個体の継続として対応づけられる時刻 $t+1$ の
個体実現を高々一つとし、
コピー個体または複数の後続実現が生じる分岐を
Phase~1 の対象から除外する。
構造的自己 $S_{\mathrm{struct}}(\Sigma_i)$
（定義~\ref{def:s_struct}）および
時間的自己同一性の予備条件
（定義~\ref{def:time_identity}）は、
この非分岐性を前提として記述されている。
組織ACSの分割・分社も、同一組織の後続が複数になる場合として、この問題に含まれる。

\textbf{問題の定式化：}
分岐（コピー個体・複数の後続実現）を許す場合について、
以下を問う。

\begin{enumerate}
  \item[(i)] \textbf{同一個体性関係の再定式化：}
  時刻 $t$ の個体実現に対して複数の後続実現が生じる場合に、
  いずれを同一個体 $\Sigma_i$ の継続とみなすか、
  複数の実現を同一系譜上の継続として扱うか、
  または分岐後の各実現を別個体として扱うかを定める
  同一個体性関係の形式化。

  \item[(ii)] \textbf{構造的自己と時間的自己同一性の拡張：}
  分岐を許す枠組みにおける構造的自己
  $S_{\mathrm{struct}}$ と時間的自己同一性の予備条件
  （定義~\ref{def:time_identity}）の再定式化、
  および崩壊 $\mathrm{Col}_t$
  （\S\ref{sec:collapse_cond}）の分岐下での扱い。

  \item[(iii)] \textbf{分岐と閉塞・崩壊の関係：}
  分岐の成立と閉塞
  $A_t=\emptyset$ および崩壊
  $\mathrm{Col}_t=1$ との論理的・動学的関係、
  ならびに分岐後の複数実現を含む個体群または系譜としての
  存続をどのように記述するかの整理。
\end{enumerate}

\textbf{現状の障壁：}
Phase~1 は非分岐性（仮定~\ref{ass:no_branching}）のもとで
単一の後続実現を前提としており、
分岐下の同一個体性は、個体の存続に関する新たな
形式的同一性関係を要する。
この関係は状態遷移写像 $f$ の一価性とは別の個体同一性上の条件であり、
Phase~1 でも非分岐性は $f$ とは別の仮定として置いている
（仮定~\ref{ass:no_branching}）。

\textbf{Phase~2以降での取り組み方針：}
個体実現間の時間的継続関係、集合値対応、または系譜構造等を
候補として、複数の後続実現を許す同一個体性の形式体系を導入し、
仮定~\ref{ass:no_branching} のもとでの非分岐的時間同一性を
その特別な場合として回収できる枠組みを検討する。
\end{problem}

\begin{problem}[T-STABパラメータの較正と予測的解釈]
\label{prob:stability}
T-STABで所与とされる
$\underline{s}_{\mathrm{SM}}$、
$\delta_{\max}$、
および$W_T$について、
対象ACSの実装種別と整合する定量的較正条件を明らかにする。
さらに、
\begin{enumerate}
  \item SM評価の参照スケールおよび下限参照値の選択に対する感度；
  \item 窓幅$W_T$と変動上限$\delta_{\max}$に対するロバスト性；
  \item T-STAB違反と閉塞・崩壊・警告条件との
  追加条件付き接続の識別可能性；
  \item 実装データからのパラメータ推定および妥当性検証
\end{enumerate}
を検討する。
Phase~1のT-STAB定義は、
これらの一般的な較正則または予測性能を与えない。
\end{problem}

\begin{problem}[推定誤差と内部更新則の定量的同定]
\label{prob:estimation_error}
TIPAのM1推定誤差、M2評価誤差、内部モデル更新、
偏向性ごとの偏向強度 $\boldsymbol{\beta}_t$（各偏向性の参照基準と尺度の指定を含む）、
$d_{\mathrm{eff},t}$、$\pi_{\tau,t}$、
四健全度等について、実装ごとの読み出し規則、担体の遷移と誤差伝播を同定し、
T-STAB、可行行動集合および閉塞到達への影響を評価する。
Phase~1の抽象的な型保存条件から、一般的な収束率や
誤差分布を導くことはできないため、追加モデルとデータを要する。
\end{problem}

\begin{problem}[核維持応答の獲得条件]
\label{prob:kernel_response_acquisition_conditions}
\textbf{背景：}
命題~\ref{prop:kernel_maintenance} は、
$x_T\in\mathrm{Viab}_f(K)$ のもとで、
すべての $s\ge T$ について $e_s\in E_{\mathrm{req}}(x_s)$ と
$u_s\in U^{\mathrm{Viab}}(x_s)$ が成り立つならば、
すべての $s\ge T$ について $x_s\in\mathrm{Viab}_f(K)$ かつ $A_s\neq\emptyset$ であることを示す。
このうち後者の前件は、実効行動の生成機構（\S\ref{sec:tipa}）と
内部モデルの更新（\S\ref{sec:extended_dynamics}）を通じて成立する場合も成立しない場合もある
（定義~\ref{def:kernel_response_acquisition}）。
本論文は、核維持応答の獲得が起こること、内部モデルの更新が獲得を導くこと、
および $U^{\mathrm{Viab}}(x)\neq\emptyset$ から実際に選ばれる行動が $U^{\mathrm{Viab}}(x)$ に属することの
いずれも一般的な帰結として置かない（備考~\ref{rem:kernel_response_scope}）。
また ACS が $\mathrm{Viab}_f(K)$ や $U^{\mathrm{Viab}}(x)$ を内部に保持または計算することも要求しない
（\S\ref{subsec:description_levels_global}）。

\textbf{問題の定式化：}
\begin{enumerate}
  \item[(i)] 観測（$P_{\mathrm{sens}}$）と内部推定が、
  核維持応答とそれ以外の応答を区別するのに要する情報を保つための条件を特定すること。
  \item[(ii)] 内部モデルの更新と応答の形成・選択の機構に関するどの追加条件のもとで、
  ある時刻以降に核維持応答が選ばれ続ける（獲得が成立する）か、
  およびその時刻の上界を与えうるかを検討すること。
  \item[(iii)] 必要環境の充足 $e_s\in E_{\mathrm{req}}(x_s)$ が途切れる外部環境状態の列のもとで、
  核維持応答の選択と、実現する軌道が $\mathrm{Viab}_f(K)$ にとどまることとの関係を整理すること。
\end{enumerate}

\textbf{現状の障壁：}
Phase~1 は可行性核と核維持応答を形式量として定義するが、
内部機構がそれらを内部情報だけからどのように近似し選択するかについては、
一般的な条件を与えない。

\textbf{取り組み方針：}
決定論的な有限の実現における条件はPhase~1.$\alpha$の詳細化として、
推定の不確実性や外部環境の確率的変動を含む一般の動学への拡張はPhase~2として検討する。
\end{problem}

\begin{problem}[関係記述から因果モデルへの拡張]
\label{prob:relational_causality}
\S\ref{sec:rel_min}の記述的関係構造を保ちながら、
必要な応用に限って、対象別の因果寄与を定義する条件を明らかにする。
具体的には、介入規則、比較軌道、反実仮想、交絡する複数関係、
時間窓の選択およびサイズ汎関数の選択が
因果的解釈へ与える影響を検討する。
記述的関係カテゴリーそのものを因果カテゴリーへ読み替えないことが前提となる。
\end{problem}

\begin{problem}[閉塞・崩壊・回復の観測可能性]
\label{prob:collapse_observability}
有限観測から閉塞・崩壊・回復可能性を判定するために、
どの観測量、履歴長、内部状態へのアクセスおよび外的作用モデルが必要かを明らかにする。
あわせて、閉塞到達時刻・崩壊到達時刻の予測、
経路集合への流入条件、回復判定のロバスト性、
誤判定の上界等を定式化する。
\end{problem}

\begin{problem}[$K$ への多段の回復可能性と必要環境要因]
\label{prob:multistep_recovery}
\textbf{背景：}
\S\ref{sec:three_principles}の一段階臨界必要要因、EUおよび局所FRは
一段の条件である（備考~\ref{rem:eu_one_step_scope}）。
$x_t\notin K$ の時刻の閉塞は、一段では $K$ へ戻れないことだけを表す。
本論文は、そこから有限段で戻れないことを一般的な帰結として置かない
（\S\ref{sec:acs_def}、備考~\ref{rem:blockage_scope_sec2}）。
また $\mathrm{Viab}_f(K)\subseteq K$ であるから、
$x\notin\mathrm{Viab}_f(K)$ は多段の回復不能の根拠にならない。

\textbf{問題の定式化：}
\begin{enumerate}
  \item[(i)] $K$ の外の状態から有限段で $K$ へ到達できることを、
  行動の列と外生背景成分の列の量化順序
  （行動の列を背景成分に先立って固定するか、各段の状態と外部環境状態に応じて選ぶか）
  を明示して定義すること。
  \item[(ii)] その到達に要する必要環境要因を、
  一段階臨界必要要因との関係を含めて特徴づけること。
  \item[(iii)] 多段の回復の途中にある時刻の閉塞と、
  有限段での回復不能とを区別して記述すること。
\end{enumerate}

\textbf{現状の障壁：}
Phase~1は多段の到達可能性を定義せず、
必要環境要因との接続も一段の条件に限っている。

\textbf{取り組み方針：}
決定論的な状態遷移と環境実現（外生背景成分の列と $g_O$）のもとで扱える部分は
Phase~1.$\alpha$の詳細化として、
外部環境の確率的変動を伴う到達可能性はPhase~2として検討する。
\end{problem}

\begin{problem}[組織ACSの動態・同一性・集約則]
\label{prob:social_dynamics}
\S\ref{sec:org_acs}の時刻別参加ACS集合と
組織ACSの記述的適用規約を基礎として、
Phase~2では、参加ACSの加入・離脱・生成・消滅、
組織境界・機能分担・関係構造の時間変化、
組織同一性の保持・復帰条件、
組織水準のEU・FRの扱いとT-STABの具体化、
および組織水準の閉塞・崩壊経路を扱う拡張を構成する。
個体量から組織量への集約則はモデル依存であり、
単一の平均・最小・加重平均を一般則として前提としない。
\end{problem}

\begin{problem}[能力の向上と有効な選択肢の縮小（能力逆説）]
\label{prob:capability_paradox}

\textbf{背景：}
本問題では、認知的処理能力（\S\ref{sec:acs_def}、備考~\ref{rem:cognitive_capacity_policy}）の向上が
内部モデルの表現可能性を広げ、
それが行動の候補（\S\ref{sec:tipa} の $\mathrm{Cand}_t$ 等）を生成しうる幅の拡大に
接続しうる、という仮説を考える。
一方、各時刻に状態遷移写像へ与えられる行動入力は $u_t\in U$ の一つであり
（定義~\ref{def:state_transition}、備考~\ref{rem:f_reading}）、
この点は能力の高低によらない。
本論文では、能力の向上に伴って、自己を維持するための管理負荷と、
他のACSその他の外部対象との関係から生じる制約とが増大し、
長い時間地平で評価したときに有効な選択肢が少数に絞られうる、
という見方を\emph{能力逆説}と呼ぶ。
さらに、時間地平の見通しが長い系ほど、
ACS同士の可行性を維持する協調的な行動が
有効な選択肢として残りうるという見方を、
能力逆説から導かれうる帰結の候補として記録する。
これらは本論文の仮説であり、Phase~1 の定義・命題から得られた一般的な帰結としては主張しない。
ここでいう管理負荷と有効性は、Phase~1 では形式的に定義していない。

\textbf{問題の定式化：}
以下を問う。

\begin{enumerate}
  \item[(i)] \textbf{潜在的な選択肢と有効な選択肢の区別：}
  内部モデルが生成しうる行動の候補（\S\ref{sec:tipa} の $\mathrm{Cand}_t$ 等）と、
  時間地平を指定したときに有効とされる行動の集合とを区別して定義すること。
  有効性の基準には、一段階の可行行動集合 $A_t$（定義~\ref{def:at}）を
  多段の時間地平へ拡張したものが候補となる。
  ただし、決定論的な状態遷移のもとで初期値、外生背景成分の列および $g_O$ を固定すると、
  $T+1$ 段にわたって可行性を保つ行動列の最初の $T$ 段は
  $T$ 段にわたって可行性を保つ行動列であるから、
  そのような行動列の最初の行動がとりうる集合は、
  時間地平を長くしても大きくならない。
  この縮小は能力の高低によらず定義から生じる。
  したがって能力逆説の内容は、この縮小の有無ではなく、
  能力の向上が縮小の程度をどのように変えるかにある。

  \item[(ii)] \textbf{能力と制約の結合：}
  能力の向上が、量的容量・時間応答性・統合の質
  （備考~\ref{rem:cognitive_capacity_policy}）の変化、
  外部対象との関係（\S\ref{sec:rel_min}）その他を通じて、
  有効な選択肢の集合に与える影響を形式化すること。
  能力の向上が可行行動集合を広げる場合や、
  制約そのものを変える場合のように、
  逆向きに働く効果も含めて扱う。

  \item[(iii)] \textbf{協調的行動との関係：}
  複数のACSが互いの外部環境の一部となる状況で、
  時間地平の見通しの長さと、
  ACS同士の可行性を維持する行動が有効な選択肢として残ることとの関係を定式化すること。
  協調的な行動が有効な選択肢に含まれうることは、それだけでは、
  協調的でない行動を有効な選択肢から除く条件を定めない。
  したがって、有効な選択肢が協調的な行動に絞られるための条件そのものを問う。
\end{enumerate}

\textbf{現状の障壁：}
Phase~1 は単一ACSの決定論的記述を中心とし、
複数のACSの同時的な時間発展、
時間地平にわたる有効性の基準、
およびACS同士の可行性の相互依存を形式化していない。
\S\ref{sec:rel_min} の関係記述は因果的寄与を含まず、
\S\ref{sec:org_acs} は組織水準への限定的な構造適用にとどまる。

\textbf{Phase~2以降での取り組み方針：}
Phase~2 では、問題~\ref{prob:social_dynamics}の拡張と合わせて、
複数のACSの相互作用と、時間地平にわたる有効な選択肢の集合を形式化する。
Phase~3 では、問題~\ref{prob:ecit_full}の中で、
有効な選択肢が協調的な行動に絞られる条件と
ECITの協調概念との関係を検討する。
いずれの段階でも能力逆説を前提として採用せず、
成り立つ条件と成り立たない条件の双方を検討の対象とする。
\end{problem}

\begin{problem}[横断的対応付けと実証的検証]
\label{prob:cross_domain}
人間・生物、AI・ロボット、組織等について、
ACSの7タプル、三原理、TIPA、PSL・ISL、
\S\ref{sec:extended_dynamics}変数、関係記述子、閉塞・崩壊の操作的対応を定め、
測定可能性、再現性、予測性能および反証可能性を評価する。
形式的な対応例と、経験的に支持された対応を明確に区別する必要がある。
\end{problem}

\begin{problem}[ECITとの正式接続]
\label{prob:ecit_full}
ACS理論で与えた可行性保存の構造と、
ECITにおける倫理的協調、愛・智慧・責任等の概念との
正式な対応関係を構成する。
その際、記述的な可行性条件から規範的価値判断が
自動的に導かれるとは仮定しない。
また、問題~\ref{prob:social_dynamics}で扱う
Phase~2の組織動態の結果については、
記述的な結果としての位置づけを維持したまま、
ECITとの接続を検討する。
\end{problem}

\begin{remark}[未解決問題の依存関係]
\label{rem:problem_dependencies}
上記の問題には実務上の依存関係があるが、
厳密な論理的必須順序を主張しない。

単一ACSについては、
時間的自己同一性、T-STABの較正、推定誤差、内部更新則、
核維持応答の獲得条件、
および多段の回復可能性の検討を基礎的課題とする。
そのうえで、一つの自然な研究順序は、
\[
  \text{単一ACS}
  \;\longrightarrow\;
  \text{関係・閉塞}
  \;\longrightarrow\;
  \text{組織・社会動態}
  \;\longrightarrow\;
  \text{ECIT接続}
\]
である。

問題~\ref{prob:capability_paradox}は、
組織・社会動態と、ECIT接続の双方にまたがる。
横断的実証検証は各段階と並行して進めうる。
\end{remark}

%----------------------------------------------------------------------
\subsection{Phase~2・3への展望}
\label{subsec:phase_outlook}

本節では、Phase~1の定義・命題を比較基準として保持しつつ、
その強い境界条件を段階的に弱化する研究方向を示す。
以下は研究計画であり、達成を保証するものではない。

\subsubsection*{Phase~2：確率的・動態的拡張}
\label{subsubsec:phase2}

\noindent\textbf{Stage~1（単一ACSのロバスト化）：}
問題~\ref{prob:branching_identity}、
問題~\ref{prob:stability}、
問題~\ref{prob:estimation_error}および
問題~\ref{prob:multistep_recovery}を中心に、
分岐を許す時間的自己同一性、
T-STABパラメータの較正、
推定誤差、更新則、到達性および安定性を
具体モデル上で検討する。
同時に、Phase~1で採用した
決定論・離散時間・コンパクト性・非分岐性等の仮定を、
どの順序で弱化できるかを評価する。

\medskip
\noindent\textbf{Stage~2（関係・閉塞・回復の動態化）：}
問題~\ref{prob:relational_causality}および
問題~\ref{prob:collapse_observability}を中心に、
関係記録と因果推定を分離したまま、
確率的・連続時間的な関係変化、閉塞到達、崩壊到達、
回復可能性および観測誤差を扱う。

\medskip
\noindent\textbf{Stage~3（組織・社会動態の拡張）：}
問題~\ref{prob:social_dynamics}のPhase~2部分を中心に、
参加ACS集合の時間変化、加入・離脱、
組織同一性、機能分担、関係構造、
組織水準のEU・FRの扱いとT-STAB、
および組織水準の閉塞・崩壊経路を
モデル依存の明示的規則として定式化する。
個体水準から組織水準への集約則は、
適用対象ごとに型と根拠を明示し、
一般則として無根拠に固定しない。
あわせて、問題~\ref{prob:capability_paradox}のうち、
複数のACSの相互作用と時間地平にわたる有効な選択肢の集合の形式化を扱う。

\medskip
\noindent\textbf{Stage~4（実証的評価）：}
問題~\ref{prob:cross_domain}を通じて、
各ドメインにおける観測可能量と理論変数の対応を検証する。
Phase~2で導入する追加構造が予測力・説明力・反証可能性を
実際に改善するかを評価し、必要に応じてPhase~1の主張範囲を再限定する。
あわせて、本論文が目指す共有可能な記述言語としての性質、
すなわち、異なる観察者が同じ主体の同じ行動を
ACSの語彙で記述したとき、
その記述がどの範囲で一致し、どこでモデル化の選択によって分かれるかを、
具体的な対象について検討する。

\subsubsection*{Phase~3：ECITとの正式接続}
\label{subsubsec:phase3}

Phase~3では、問題~\ref{prob:ecit_full}を中心に、
ACS理論の記述的な可行性保存構造と
ECITの倫理的協調概念との正式な接続を検討する。
問題~\ref{prob:capability_paradox}のうち協調的な行動に関わる部分も、
この段階で扱う。
この接続においてACS理論が担う役割は、
協調の当事者となる主体それぞれの行動を、
特定の社会や文化の規範を前提とせずに、
共通の構造のもとで記述する層を与えることである
（\S\ref{subsec:background}）。
ECITが扱う協調は、
異なる規範のもとにある主体どうし、
あるいは人間と自律的AIシステムのように成り立ちの異なる主体どうしの間で
問題となる。
そうした主体の間では、
ある主体がなぜそう行動したのかを、
いずれか一方の規範によって判定する前に、
どの可行性上の要請が関わり、
どの内部機構の作動を経てその行動が成立したかを
共有可能な形式で表現できることが、
協調を論じるための前提の一つになると考えられる。
Phase~3では、この記述層の上に
ECITの協調概念をどのように構成できるかを検討する。

ここでは、ACS理論が与える構造的記述、
ECITが扱う協調の条件、
および行動の善し悪しに関する価値判断を、
互いに異なる論理的地位をもつものとして区別する。
Phase~1で定義したACS構造から、
特定の価値判断または社会規範が論理的に自動導出されるとはしない。
共通の記述層を与えることと、
その記述にもとづいて行動の善し悪しを定めることとは、
別の論理的地位をもつ。
規範的要素を導入する場合には、その追加公理・評価基準・
社会的前提をACSの構造的定義と明確に分離する。
また、社会規範そのものを判定の基準としてではなく、
主体が置かれた環境、関係、または内部モデルの内容として
記述の対象に含める可能性も、
この段階で検討の対象となる。
Phase~1は、規範をこのように記述する形式を与えていない。

\begin{remark}[Phase区分の研究戦略的意義]
\label{rem:phase_strategy}
Phase~1--3の区分は、研究の時系列だけでなく、
異なる証明強度とモデル依存性を分離するための区分である。
Phase~1は決定論的・離散時間の最小構造、
Phase~2は確率・連続時間・因果・組織／社会動態等の拡張、
Phase~3はECITとの正式接続を担当する。
後続Phaseの結果によりPhase~1の適用範囲が
修正または限定される可能性をあらかじめ排除しない。
\end{remark}

%----------------------------------------------------------------------
\subsection{結論}
\label{subsec:conclusion}

本論文は、自律認知システム理論（ACS理論）のPhase~1として、
自律的認知システムを
$\Sigma=(X,U,E,f,\Pi,M,\mathrm{Update})$という7タプルで記述し、
可行性保存を中心とする構造的枠組みを構成した。
以下では、本論文が確立した範囲と、確立していない範囲を分けて確認する。

\subsubsection*{Phase~1で確立した主要構造}
\label{subsubsec:core_results}

\begin{itemize}
  \item \textbf{ACSの形式的定義：}
  \S\ref{sec:acs_def}で、状態・行動・必要環境・状態遷移・持続評価・内部モデル・更新作用素からなる
  7タプルと、可行集合$K$、可行行動集合$A_t$その他の基本構造を定めた。
  外部環境状態空間を外生背景成分と行動応答成分に分け、
  構造的自己と整形式実現の区別のもとで、機能的終端と崩壊を状態条件として定義した。

  \item \textbf{三原理：}
  \S\ref{sec:three_principles}で、SM・EU・FRを
  同一の論理地位へ平板化せず、それぞれの時刻別充足述語・定義域・証明強度を区別して形式化し、
  閉塞を臨界必要要因の充足の有無によって二分した。
  また、SM評価・EU・局所FRを有限時間窓上で統合する
  三原理時間統合評価条件（T-STAB）を定義した。

  \item \textbf{TIPA・PSL・ISL：}
  \S\ref{sec:tipa}で、M1情報統合、M2適応的評価、M3応答行動からなるTIPAと、
  その物理的基盤層PSL、実装部位層ISLの機能的構造を整理した。

  \item \textbf{時刻間動学：}
  \S\ref{sec:extended_dynamics}で、原始的な機能の標準写像と読み出し規則を固定し
  環境実現を与えたとき、一段の作動記録がただ一つ定まり次の時刻へ閉じることを示し、
  記憶、偏向性（参照基準に相対する偏りの傾向）とその程度である偏向強度、
  実効的発展自由度、予測精度、実効的時間地平、健全度等を、
  状態・内部モデル・外部環境状態または作動記録の有限区間からの読み出しとして時刻$t$から$t+1$へ接続した。

  \item \textbf{関係性の最小構造：}
  \S\ref{sec:rel_min}で、外部対象との単方向・双方向関係、方向性、評価規則に相対する関係品質、代表カテゴリーおよび六系分類を、
  因果帰属を自動的には含まない記述・診断・比較構造として整理した。

  \item \textbf{閉塞・崩壊：}
  \S\ref{sec:collapse_cond}で、閉塞$A_t=\emptyset$と
  崩壊$\mathrm{Col}_t=1$を明確に分離し、
  外的回復可能性、軌道記録、閉塞経路集合、崩壊到達経路集合および到達時刻を整理した。

  \item \textbf{組織ACSへの限定的構造適用：}
  \S\ref{sec:org_acs}で、ACS集団、ACS組織および組織ACSを区別し、
  組織自身について整形式な7タプル構造が成立する場合に限って、
  個体ACSの参加の有無によらず組織ACSとして扱う記述的適用規約を与え、
  参加ACS数により個体ACS参加型と個体ACS非参加型を区別した。
  組織水準のACS構造が抽象構造型の実現であることにより、
  先行節の三原理、閉塞・崩壊および経路集合の定義を組織水準へ適用する形を与え、
  組織水準の可行集合・可行行動集合と閉塞の特徴づけを示した。
  個体量から組織量への普遍的集約則と、
  組織についての判定規則の一般形式は与えなかった。
\end{itemize}

\subsubsection*{理論の位置づけ}
\label{subsubsec:positioning}

ACS理論は、価値判断そのものを与える規範理論ではなく、
自律認知システムの可行性保存に関わる条件を記述する構造理論として位置づけられる。
本論文が既存のViability Theory、凸解析、動的システム論、制御理論、
認知アーキテクチャその他の理論を包含または一般化したとは主張しない。
本論文で述べる既存理論との関係は、明示した部分における整合性・接続可能性に限定される。

本論文が構造理論として記述を目指したのは、
自律的に振る舞う主体の行動を、
善悪の判定に先立って、
可行性上の要請と内部機構の作動との関係として記述するための
共通の形式である。
Phase~1で与えたのはその決定論的・構造的な基礎であり、
この形式が異なる領域や文化のもとにある主体の行動を
実際にどの程度共有可能な形で記述できるかは、
問題~\ref{prob:cross_domain}および
問題~\ref{prob:ecit_full}として残る。

また、Phase~1の形式的結果が人間、生物、AI、ロボット、組織に
経験的に妥当であることは本論文だけからは導かれない。
その評価は問題~\ref{prob:cross_domain}として残る。

\subsubsection*{主張強度の原則}
\label{subsubsec:principle_check}

本論文では、各結果の論理的地位を
定義的帰結、仮定付き命題、傾向的記述、
未解決問題等に分け、次の原則を維持する：
\[
  \text{主張の強さ}\;\leq\;\text{証明の強さ}
\]
証明または定義から導かれない内容については、
一般的含意、因果効果、安定性、最適性、実証的妥当性または規範的正当性を
自動的に付与しない。

本論文で提示したPhase~1の枠組みが、
自律認知システムと倫理的協調知性の将来的な数理研究に
どの程度有効な基盤となるかは、Phase~2以降の形式的・実証的検証によって評価される。

\medskip
本論文およびACS理論は平和利用を志向して公開される。
ACSによる協調的な社会構造、ならびに人間と他の知性体の
持続的な共生のあり方の探求に資することは、
Phase~1から証明された帰結ではなく、研究上の理念的な希望として位置づける。

\begin{remark}[AI支援の位置づけ]
\label{rem:ai_assistance}
本研究では、言語モデルを、原稿案・数式・証明候補・検証用コードの作成の補助、
文献候補の整理、および論理的整合性・記号整合性・参照整合性の確認に利用した。
理論上の採否、修正判断および最終的な研究責任は著者に帰属する。
著者以外の独立した専門家による本文と検証用コードの全面的な検証は行われていない。
\end{remark}
